% !PRINTED 386
\weberpage{386}
%<<CONT
when $a$ is a \gesperrt{quadratic non-residue} of $l$, and by the extended definition
of this symbol introduced by \gesperrt{Jacobi},
\weq{2}{\left(\frac{a}{l\,l'\,l''\dots}\right) = \left(\frac{a}{l}\right)\left(\frac{a}{l'}\right)\left(\frac{a}{l''}\right)\cdots}

If furthermore $a$ is odd, then
\weq{3}{\left(\frac{-1}{a}\right) = (-1)^{\frac{a-1}{2}}, \quad \left(\frac{-2}{a}\right) = (-1)^{\frac{a^2-1}{8}},}
and from this, if $a, b, \dots$ are relatively prime to $m$:
\weq{4}{\left(\frac{a}{m}\right)\left(\frac{b}{m}\right)\cdots = \left(\frac{ab\dots}{m}\right);}
if $D$ is a discriminant and $a$ a number coprime to $2D$ which is
properly representable by a form $\varphi$ of the discriminant $D$,
 and if further $l$ is an odd prime dividing $D$,
then the value of the symbol $\left(\frac{a}{l}\right)$ depends not on the particular
number $a$, but only on the class of forms $A$ to which $\varphi$ belongs.
The same holds for
\weq{5}{
\begin{gathered}
\left(\frac{-1}{a}\right) \quad D \equiv 0,-4 \pmod{16},\\
\left.\begin{aligned}
\left(\frac{2}{a}\right)& \quad D \equiv 0,8\\
\left(\frac{-2}{a}\right)& \quad D \equiv 0,-8
\end{aligned}\right\} \pmod{32}.
\end{gathered}
}

From this the characters of the classes of forms can be
composed.

In a later application we shall have to take special account
 of the case in which $D$ is divisible by 4.
We therefore set:
\weq{6}{D = -4m\footnotemark,}
\footnotetext{In the notation of \gesperrt{Gauss: forms of the first kind of the determinant $-m$.}}
and for this case we obtain, if $l, l', \dots$ are the different
odd primes dividing $m$, the following
characters:
%CONT>>
